\section{Background: Semantic Table Interpretation}
What is it....\todo{improve}

A vast amount of information is provided as structured data on the Web in tables, this quantity grew over the years. This increase can be linked to the uptake of the Open Data movement, whose purpose is to make a large number of tabular data sources freely available, addressing a wide range of domains, such as finance, mobility, tourism, sports, or cultural heritage~\cite{Neumaier2016}. The massive availability of tabular data on the Web makes Web tables a valuable source to consider for data miners. For instance, these tables can be employed for data integration tasks or to construct and extend \acp{kg}~\cite{Kejriwal2021}. In this field, the table-to-KG matching problem, also referred to as \ac{sti}, is the process of adding the semantic meaning of a table by mapping its elements (i.e., cells/mentions, columns, rows) to semantic tags (i.e., entities, classes, properties) from \acp{kg} (e.g., Wikidata, DBpedia). This process is typically broken down into the following tasks: (i) cell/mentions to \ac{kg} entity matching (CEA task), (ii) column to \ac{kg} class matching (CTA task), and (iii) column pair to \ac{kg} property matching (CPA task)~\cite{Cutrona2021-1}. In the last decade, the table-to-KG has collected much attention in the research community~\cite{Cutrona2021-1}. This interest is also certified by the introduction of the ``SemTab: Semantic Web Challenge on Tabular Data to Knowledge Graph Matching'' international challenge\footnote{\href{https://www.cs.ox.ac.uk/isg/challenges/sem-tab/}{www.cs.ox.ac.uk/isg/challenges/sem-tab/}}, now in its 4th version. The challenge consists of different rounds in which groups of tables with different features and levels of difficulty have to be annotated.